\chapter{Agents and tool use}
\label{ch:agents}

\section*{What this chapter is for}

\reported{Agent and tool-use questions are very common and the category is
growing, and agentic exercises are the second-largest group in analysed
take-home assessments.}{field-guide} The questions are less standardised than
retrieval's, which makes the method of Chapter~\ref{ch:answering} matter more
here than anywhere else.

\section*{The questions}

\begin{bankq}{When is an agent the wrong solution?}
\qasked
The counterpart of retrieval's opening question, and asked for the same reason.

\qtested
Whether you default to the architecture that is currently fashionable.
\reported{A candidate is recorded as receiving a negative assessment for
repeatedly proposing model-based solutions to a problem that did not need
one.}{field-guide-trends}

\qstaff
An agent is a loop that chooses its own next step. You are buying the ability
to handle inputs you did not enumerate, and paying for it in non-determinism,
latency and a much larger surface of things that can go wrong.

So it is wrong whenever the steps are known. A fixed sequence of three calls is
a function; implementing it as an agent means the model can decide not to make
one of them, and now you have a reliability problem you constructed yourself.
It is wrong when every path must be auditable in advance. It is wrong when
latency is bounded tightly, because a loop has no upper bound you control
without imposing one. And it is wrong when the cost of a mistaken action is
high and the loop can act without a person \dash{} unless the boundary stops
it, which is the next question.

\judgement{The useful reframing is that an agent is a control-flow choice, not
a capability. Asking ``would I write this as a loop with a dynamic branch, and
would I let the branch be decided by something I cannot test exhaustively?''
turns it into a question you already know how to answer.}

\qsenior
Naming cost and latency as the trade-off, correctly, and stopping there. It
treats the decision as a performance question when it is a determinism
question.

\qcontested
How much agency is useful is the live argument in this field, and the answer
has moved more than once. \judgement{The durable part is not where the line is
but that it is a line at all: anybody who says agents are always or never right
is describing a preference.}

\qdrill
Take a workflow you have automated conventionally and say what making it
agentic would buy and cost. If the answer is that it buys nothing, say why
somebody might still propose it \dash{} that second half is the interesting one.
\end{bankq}

\begin{bankq}{How do you stop an agent doing something irreversible?}
\qasked
Sometimes as a design question, sometimes as a scenario about an action that
should not have happened.

\qtested
Where you put the control. This question separates people who have operated one
of these from people who have built one, in a single answer.

\qstaff
\textbf{The tool boundary is the security boundary; the agent's behaviour is
user experience.} An instruction telling the model to ask before acting is a
request, and it is subject to everything that can influence the model
\dash{} including content that arrives inside the data the system itself
fetched.

So the confirmation lives in the tool. The write tool refuses to act without a
token; the token is issued only when a person approves a specific action; it is
bound to exactly what was shown to them, and it is single-use. An agent that is
argued into calling the tool early then fails at the tool, which is a mechanism
rather than a hope.

\judgement{The design test is to ask what happens if the model does exactly the
wrong thing on purpose. If the answer depends on the prompt, there is no
control. If the answer is that the call is refused, there is.}

That reasoning extends: least privilege per tool rather than one capable
credential; a read tool that cannot write even by accident; and a blast radius
you can state \dash{} the worst thing this loop can do without a person is a
sentence you should be able to say about any system you have built.

\qsenior
A correct description of a confirmation step in the agent's flow, with the
prompt instructing it to ask first. It works when nothing is adversarial and
its failure mode is silent.

\qcontested
How much to enforce at the boundary versus in the loop is a real engineering
trade-off, because boundary enforcement costs flexibility.
\judgement{For irreversible actions there is no trade-off worth making. For
reversible ones the loop is fine, and being able to say which class an action
is in is most of the answer.}

\qdrill
Take a tool you would give an agent and write its refusal condition: the
precise circumstance in which it declines, expressed without reference to the
model's intentions.
\end{bankq}

\begin{bankq}{How would you design the tools themselves?}
\qasked
Often the most open question in the round, and the one where the method from
Chapter~\ref{ch:answering} does the most work.

\qtested
Whether you can design an interface for a caller that is capable, literal and
occasionally confidently wrong. This is API design with an unusual consumer,
which is why it transfers so well from backend work.

\qstaff
Few tools, each doing one thing, named for the outcome rather than the
implementation. A model choosing between fifteen overlapping tools chooses
badly, and the fix is fewer tools rather than a longer description.

Parameters that are hard to get wrong: enumerations over free strings,
identifiers that came from a previous result rather than being constructed, and
required fields that genuinely are required. \judgement{The single highest-value
constraint is that an identifier passed to a write must have appeared in an
earlier result in the same session. It makes a confidently invented identifier
structurally impossible rather than something you hope to catch.}

Errors that say what to do next. A tool returning ``invalid request'' gives the
model nothing to act on; one returning ``that identifier was not found; call
\code{list\_items} first'' produces a recovery. You are writing error messages
for a reader who will act on them literally.

And a description that says when \emph{not} to use it, which is the part
everybody omits and the part that resolves most wrong-tool selection.

\qsenior
Clear names, good descriptions, typed parameters. Correct, and it is the
answer for a human consumer. The literal-caller constraints are what makes it
this domain.

\qcontested
Granularity is genuinely contested \dash{} many small tools against few
composite ones \dash{} and it interacts with the model and the task in ways
nobody has characterised well. \judgement{Start coarse. Splitting a tool is
cheap; unpicking a model's habit of calling three tools where one would do is
not.}

\qdrill
Take an API you have designed and rewrite three endpoints as tools. Then write
the sentence saying when each should not be used. That sentence is usually the
hardest one and it is the one that reveals whether the tool has a clear purpose.
\end{bankq}

\begin{bankq}{How do you know the agent is still behaving after a change?}
\qasked
Usually arrives as a question about testing, and is the bridge to
Chapter~\ref{ch:evals}.

\qtested
Whether you test the behaviour or the text. The answer given here is what makes
the evals chapter land.

\qstaff
Test what it \emph{did}, not what it said. The sequence of tool calls, their
arguments, and the events emitted is a structure you can assert against; the
reply text is prose that changes with every prompt edit and turns a rewording
into a false regression.

That reframing brings a second and less obvious point: half the assertions
should be about what did \emph{not} happen. An agent that refuses politely and
calls the tool anyway has passed a check that only looked for the refusal. Any
scenario about declining to act needs both halves \dash{} the refusal, and the
absence of the call.

\qsenior
A set of cases with expected outputs compared by similarity or by a model-based
grader. It works, it is expensive, and it is unstable against changes that
altered nothing important.

\qcontested
Whether behaviour testing is sufficient for systems whose output quality is the
product. \judgement{It is not sufficient and it is the part that can be
mechanical and cheap. The right structure is a deterministic layer that gates
every change and a slower judged layer that does not \dash{}
Chapter~\ref{ch:evals}.}

\qdrill
For an agent you know, write three assertions about a single turn that do not
mention any text it produced. If you cannot get to three, the system is not
emitting enough about what it did, which is an observability finding rather
than a testing one.
\end{bankq}

\section*{What this chapter does not claim}

The frequency and the negative-assessment anecdote come from one corpus. The
anecdote is one interviewer's report of one candidate and is used to illustrate
a position, not to establish one.

The tool-boundary argument, the identifier-grounding constraint and the
start-coarse recommendation are the author's judgement. The first is close to
consensus among people who have operated these systems; the other two are
positions.

No protocol is named anywhere in this chapter. Protocol-specific questions
appear in content-marketing material far more often than in any primary account
of an interview, and Appendix~\ref{app:refused} says so.
